\chapter{Beriberi}
\index{Beriberi}

\section*{Definition and Overview}
Beriberi is a disease caused by a deficiency of thiamine (vitamin B1), an essential nutrient required for converting food into energy. There are two principal forms: \textbf{dry beriberi}, which mainly affects the peripheral nerves, and \textbf{wet beriberi}, which chiefly affects the heart and circulation. A related, life-threatening neurological emergency called Wernicke encephalopathy can progress to Korsakoff syndrome, a condition of permanent memory impairment.

Beriberi is readily treatable with thiamine replacement if it is recognised in time, and it is entirely preventable through an adequate diet. It is historically associated with diets based on polished white rice, but in developed countries it now occurs mainly in people with alcohol dependence and in other vulnerable groups with poor nutrition.

\section*{Epidemiology}
Beriberi was historically common in parts of Asia where polished (white) rice formed the dietary staple, because milling removes the thiamine-rich outer layers of the grain. It is now rare in developed countries but persists among specific at-risk groups.

\begin{itemize}
    \item \textbf{Alcohol dependence:} in many countries and other developed countries, beriberi is seen most often in heavy drinkers, whose diets are poor and whose bodies absorb, store, and utilise thiamine abnormally
    \item \textbf{Poor nutrition:} prolonged vomiting, severe malnutrition, restrictive or fad diets, and eating disorders can all cause deficiency
    \item \textbf{Surgical causes:} thiamine deficiency may develop after weight-loss (bariatric) surgery because of reduced absorption
    \item \textbf{Increased requirements:} pregnant and breastfeeding women, growing infants, and people with chronic infection or fever have higher thiamine needs
    \item \textbf{Other vulnerable groups:} older people, refugees, and those subsisting on a very limited variety of foods are at risk
    \item \textbf{Endemic areas:} in some regions of Asia, beriberi still occurs in communities that rely on polished rice
\end{itemize}

\section*{Aetiology and Pathophysiology}
Thiamine is a cofactor for enzymes essential to carbohydrate metabolism and energy production. Tissues with the highest energy demands, notably the heart, peripheral nerves, and brain, depend on thiamine, so they are the first to malfunction when it is lacking.

Alcohol produces deficiency through several mechanisms simultaneously: it displaces food from the diet, impairs absorption of thiamine across the gut, reduces hepatic storage, and interferes with the activation of thiamine into its biologically active form. Deficiency therefore develops rapidly in heavy drinkers. The consequences are threefold: the heart muscle weakens and the circulation fails (wet beriberi); peripheral nerves lose their myelin sheath and cease to transmit signals normally (dry beriberi); and characteristic regions of the brain, including the thalami and mamillary bodies, are damaged, producing the confusion and eye-movement signs of Wernicke encephalopathy.

\section*{Clinical Features}
Symptoms develop gradually in chronic deficiency but may appear rapidly in acute forms of the disease.

\begin{itemize}
    \item \textbf{Dry beriberi:} symmetrical tingling, burning, or numbness in the hands and feet; muscle weakness and wasting; calf cramps; difficulty walking with a high-stepping gait; and loss of deep tendon reflexes
    \item \textbf{Wet beriberi:} shortness of breath, particularly on exertion; swelling of the legs and ankles; a rapid or bounding heart rate; fatigue; and features of heart failure such as breathlessness when lying flat
    \item \textbf{Wernicke encephalopathy:} acute confusion and drowsiness; unsteadiness and difficulty with balance; and abnormal eye movements, typically nystagmus
    \item \textbf{Infantile beriberi:} seen in breastfed infants of thiamine-deficient mothers, with rapid breathing, vomiting, heart failure, a hoarse cry, and sometimes seizures
\end{itemize}

\section*{Differential Diagnosis}
The features of beriberi overlap considerably with other conditions, which must be considered in diagnosis.

\begin{itemize}
    \item \textbf{Peripheral neuropathies:} diabetes, vitamin B12 deficiency, alcohol-related nerve damage, renal failure, and Guillain--Barr\'e syndrome
    \item \textbf{Heart failure from other causes:} coronary artery disease, hypertension, valvular disease, and viral or other cardiomyopathies
    \item \textbf{Acute confusion:} infection, stroke, head injury, hypoglycaemia, and drug or alcohol toxicity
    \item \textbf{Muscle weakness and fatigue:} anaemia, thyroid disease, and other nutritional deficiencies
    \item \textbf{Eye-movement disorders:} brainstem stroke and other neurological conditions that cause nystagmus or ophthalmoplegia
\end{itemize}

\section*{When to Seek Medical Care}
See a doctor promptly for unexplained tingling or weakness in the hands or feet, leg swelling with shortness of breath, unsteadiness, or any change in memory or concentration, particularly if you drink heavily or have a poor diet.

\textbf{Contact your local emergency services for:}
\begin{itemize}
    \item Sudden confusion, drowsiness, or loss of consciousness
    \item Severe shortness of breath, chest pain, or collapse
    \item Inability to stand or sudden difficulty walking
    \item New seizures, or abnormal eye movements occurring with confusion
\end{itemize}

Wernicke encephalopathy is a medical emergency and must be treated with thiamine urgently, before or alongside any administration of glucose.

\section*{Diagnosis and Investigations}
\begin{itemize}
    \item \textbf{History:} careful inquiry into diet, alcohol use, recent illness, vomiting, and surgery, together with a full description of symptoms
    \item \textbf{Physical examination:} assessment of the heart, lungs, and legs, plus examination of reflexes, coordination, eye movements, and gait
    \item \textbf{Blood tests:} thiamine levels can be measured directly, but treatment is frequently started on the basis of the clinical picture alone, because testing should not delay therapy
    \item \textbf{Cardiac tests:} electrocardiogram (ECG) and echocardiogram to assess heart size and function in suspected wet beriberi
    \item \textbf{Nerve conduction studies:} used to characterise peripheral nerve damage in dry beriberi
    \item \textbf{Response to treatment:} rapid clinical improvement with thiamine supports the diagnosis
\end{itemize}

\section*{Management}
\begin{itemize}
    \item \textbf{Thiamine replacement:} given by intravenous or intramuscular injection in hospital for severe or acute forms, particularly Wernicke encephalopathy, where it must be given urgently and before any glucose infusion
    \item \textbf{Oral thiamine supplements:} appropriate for mild deficiency or as follow-up after initial parenteral therapy
    \item \textbf{Treat the underlying cause:} reducing or stopping alcohol, correcting dietary deficiencies, and addressing any illness that causes poor absorption or increased requirements
    \item \textbf{Supportive care:} treatment of heart failure in wet beriberi with rest, fluid management, oxygen, and cardiac medications as required
    \item \textbf{Nutritional rehabilitation:} a balanced diet with multiple vitamins, since isolated thiamine deficiency is often accompanied by other deficiencies
    \item \textbf{Physiotherapy:} to help regain strength, balance, and walking ability after nerve damage
\end{itemize}

\section*{Complications}
\begin{itemize}
    \item Irreversible peripheral nerve damage with persistent weakness, numbness, or difficulty walking
    \item Permanent heart muscle damage with chronic heart failure
    \item Progression from Wernicke encephalopathy to Korsakoff syndrome, with severe, permanent impairment of memory and the tendency to confabulate
    \item Coma and death if wet beriberi or Wernicke encephalopathy goes untreated
    \item Falls and injuries resulting from unsteadiness and muscle weakness
    \item Long-term occupational and social disability from persistent neurological or cardiac damage
\end{itemize}

\section*{Prognosis and Outcomes}
Beriberi responds very well to thiamine treatment when it is diagnosed and treated early.

\begin{itemize}
    \item Wet beriberi often improves rapidly with thiamine and cardiac support, sometimes within days
    \item Dry beriberi improves more slowly, and some nerve damage may be permanent
    \item Early treatment of Wernicke encephalopathy usually reverses the confusion and eye signs, but delay increases the risk of permanent memory loss
    \item Korsakoff syndrome, once established, is frequently permanent and requires long-term care
    \item Outcome depends on prompt recognition, the severity of deficiency, and whether the underlying cause is corrected
\end{itemize}

\section*{Prevention and Self-Care}
Beriberi is entirely preventable through an adequate diet and attention to risk factors.

\begin{itemize}
    \item Eat thiamine-rich foods regularly, including whole grains, legumes, nuts, seeds, meat, fish, eggs, and fortified cereals
    \item Avoid relying on polished white rice as the sole dietary staple; choose wholegrain or enriched varieties where possible
    \item Limit alcohol intake, recognising that heavy drinking is the most common cause of thiamine deficiency in developed countries
    \item Take thiamine supplements if advised, for example during pregnancy, with heavy alcohol use, in chronic illness, or after bariatric surgery
    \item If you drink heavily, seek help with reducing alcohol intake and improving nutrition
    \item Seek early medical review for any symptoms of numbness, weakness, swelling, or confusion, especially if you are at risk
\end{itemize}

\section*{Sources and Further Reading}
The following organisations provide reliable information on beriberi and thiamine deficiency for patients, families, and health professionals.

\begin{itemize}
    \item World Health Organization. \textit{Resources on micronutrient deficiencies and their prevention.} https://www.who.int/
    \item MedlinePlus. \textit{Health information on beriberi and vitamin B1 deficiency.} https://medlineplus.gov/
    \item NHS. \textit{Information on vitamins and minerals, including vitamin B1.} https://www.nhs.uk/
    \item MSD Manual. \textit{Professional information on thiamine deficiency and beriberi.} https://www.msdmanuals.com/
    \item Centers for Disease Control and Prevention. \textit{Resources on nutrition and deficiency disorders.} https://www.cdc.gov/
\end{itemize}
